\section*{Chapter \ref{ch:ml:generative-models-and-synthetic-data} --- Generative Models and Synthetic Data}

\begin{solution}{exo:ml:generative-models-and-synthetic-data:1}
Under the null the held-out accuracy on 4\,000 windows has standard error $\sqrt{0.25/4\,000} = 0.0079$: 0.62 is 15
standard errors above 0.5, a clear difference. 0.51 is 1.3 standard errors: no evidence of a difference, which is not
evidence of equality; a stronger classifier or more samples might find one.
\end{solution}

\begin{solution}{exo:ml:generative-models-and-synthetic-data:2}
The leverage effect and gain-loss asymmetry: with symmetric shocks and a variance that responds to squared shocks, negative
and positive returns raise volatility equally and the distribution of returns is symmetric. A GJR or EGARCH term and skewed
shocks add them.
\end{solution}

\begin{solution}{exo:ml:generative-models-and-synthetic-data:3}
A shock to the variance decays like $(\alpha + \beta)^k$; the half-life is $\ln 0.5/\ln 0.964 = 19.1$ days.
\end{solution}

\begin{solution}{exo:ml:generative-models-and-synthetic-data:4}
The effect links a day's return to the previous twenty days' sum; inside a 20-day block the ordering of real days is kept,
so a window's last day and its preceding days are often from the same block, and the relation survives there. Across block
boundaries it is broken. Blocks of a few days destroy it; blocks much longer than twenty days keep it almost intact, and
copy more of the history.
\end{solution}

\begin{solution}{exo:ml:generative-models-and-synthetic-data:5}
Neighbouring windows share 31 of their 32 days. With a random split, almost every held-out window has a near-copy in the
training half with the same label, so the classifier recognises the period, not the distribution: 0.83 for two halves of
the same process. Splitting by year leaves no neighbours across the split.
\end{solution}

\begin{solution}{exo:ml:generative-models-and-synthetic-data:6}
The test is inside the GAN's world: a strategy trained and tested on generated paths measures the generator's regularities,
including artefacts. Here a GAN gives 10.7 inside its world and 0.52 on real data. Passing stylised-fact tests says nothing
about predictability. Test the strategy on real data it was not trained on, and compare with training on real.
\end{solution}

\begin{solution}{exo:ml:generative-models-and-synthetic-data:7}
After 1\,000 steps: excess kurtosis 2.53, two-sample accuracy 0.884, TSTS Sharpe ratio 11.2. After 8\,000: 2.49, 0.734 and
9.9 (4\,000 steps gave 1.98, 0.749 and 10.7). Longer training makes the windows harder to tell apart on average but does
not restore the tails, which a discriminator on standardised windows rarely sees, and the invented predictability stays.
The cures are architectural (heavy-tailed noise, a volatility channel as in Quant GANs) or a likelihood-based model.
\end{solution}

\begin{solution}{exo:ml:generative-models-and-synthetic-data:8}
$\log p(x) = \log\int p(x\mid z)p(z)\,dz = \log\E_q[p(x\mid z)p(z)/q(z\mid x)]\ge\E_q[\log p(x\mid z)] + \E_q[\log p(z) -
\log q(z\mid x)]$ by Jensen's inequality, and the second term is $-\mathrm{KL}(q\,\|\,p)$. For $q = N(m, \mathrm{diag}\,
s^2)$ and $p = N(0, I)$: $\mathrm{KL} = \tfrac12\sum_j(m_j^2 + s_j^2 - 1 - \log s_j^2)$, the term in the code.
\end{solution}

\begin{solution}{pb:ml:generative-models-and-synthetic-data:1}
\textbf{Part I.}
\begin{enumerate}
\item Fat tails, volatility clustering, a leverage effect, six crashes and a 20-day momentum effect. The bootstrap keeps
what fits in a block; GARCH-t keeps tails and clustering; the networks can represent any of it in principle.
\item $\alpha = 0.079$, $\beta = 0.886$, 3.6 degrees of freedom; it misses the leverage effect, the skew and the momentum.
\item \omterm{def:ml:neural-networks-for-noisy-tabular-data:mlp}{Multilayer perceptrons} on standardised 32-day windows: a VAE with a four-dimensional code and a GAN with 16-dimensional
noise and two layers of 128 units, 4\,000 \omterm{def:ml:neural-networks-for-noisy-tabular-data:adam}{Adam} steps each on one core.
\item Start from Gaussian noise and apply 100 learned denoising steps, each removing the predicted noise and adding a little
fresh noise.
\end{enumerate}
\textbf{Part II.}
\begin{enumerate}[resume]
\item Bootstrap 6, GARCH-t 4, VAE 6, GAN 3, diffusion 4 (real 6).
\item 0.491, 0.539, 0.619, 0.749, 0.693 (real against real 0.545). Overlapping windows leak across a random split: 0.83 for
real against real.
\item IC and Sharpe: real 0.078 and 0.86; bootstrap 0.062 and 0.71; GARCH-t $-0.019$ and $-0.10$; VAE 0.040 and 0.76; GAN
0.052 and 0.52; diffusion 0.030 and 0.23.
\item Bootstrap windows are close to training windows 14.5\% of the time (copies), GAN windows 11.2\% (collapse onto calm
windows), the others near the 5\% of real test windows.
\end{enumerate}
\textbf{Part III.}
\begin{enumerate}[resume]
\item 10.7, from linear autocorrelation and smooth volatility patterns the GAN invented.
\item GARCH-t has no predictability, so the fitted forecast is a random combination of features whose real-data IC is
random in sign: $-0.10$ on average with a standard deviation of 0.65 across five samples.
\item A stress test needs tails and clustering: bootstrap or GARCH-t, with the crashes; augmentation needs the
predictability without artefacts: none of the networks here is good enough, and the bootstrap's copies add little new.
\item It reproduces real 20-day sequences: anyone with the real data can recognise them.
\end{enumerate}
\textbf{Part IV.}
\begin{enumerate}[resume]
\item \emph{Paths that never happened.} Facts passed and two-sample accuracy: bootstrap 6 and 0.491, GARCH-t 4 and 0.539,
VAE 6 and 0.619, GAN 3 and 0.749, diffusion 4 and 0.693. TSTR Sharpe against 0.86 trained on real: 0.71, $-0.10$, 0.76,
0.52 and 0.23; the GAN's own world gives 10.7.
\item TSTR close to train-on-real across several samples and fits, on a test period after the training data.
\item Plant known effects in generated data, run the backtester, and check that it recovers them and nothing else.
\item When the mechanism is the question (market impact, queue dynamics, a policy change), and when data are too few to
learn a distribution.
\item Feed the regime as an input to the generator; the new risk is regimes that are rare in training and poorly generated
exactly where stress tests need them.
\item The generator's version and training data window, its seeds, the tests it passed, its purpose, and which results used
it.
\item Tail losses from diffusion samples would be understated; stress with the tails from the real data or a heavy-tailed
model.
\item Its shape, and only as much of its predictability as the architecture and the data force it to learn.
\end{enumerate}
\end{solution}

\begin{solution}{iq:ml:generative-models-and-synthetic-data:1}
Heavy tails, no linear autocorrelation, volatility clustering, the leverage effect, gain-loss asymmetry, aggregational
Gaussianity, among others. GARCH with fat-tailed shocks reproduces tails, clustering and aggregation; asymmetric variants
add leverage.
\iqlookfor{four or more facts, and which a GARCH reproduces.}
\end{solution}

\begin{solution}{iq:ml:generative-models-and-synthetic-data:2}
VAE: an encoder to a latent distribution and a decoder, trained on a likelihood bound; samples are blurry or too Gaussian.
GAN: a generator trained against a discriminator; \omterm{def:ml:generative-models-and-synthetic-data:gan}{mode collapse} and unstable training. Diffusion: learned denoising of
gradually noised data; slow sampling, and tails depend on training. Each fails differently on fat-tailed data.
\iqlookfor{the training principle and a failure for each.}
\end{solution}

\begin{solution}{iq:ml:generative-models-and-synthetic-data:3}
By purpose: stylised facts, a two-sample test with a time-blocked split, TSTR against train-on-real on later data, and a
\omterm{def:ml:large-language-models-in-finance:memo}{memorisation} check; never by results inside the synthetic data alone.
\iqlookfor{TSTR and a blocked two-sample test, tied to the purpose.}
\end{solution}

\begin{solution}{iq:ml:generative-models-and-synthetic-data:4}
Whether the improvement holds on real data not used to fit the generator; whether the generator was trained on the test
period; whether synthetic paths carry artefacts the model learned (TSTS against TSTR); and whether the same effect appears
with a bootstrap.
\iqlookfor{real out-of-sample evaluation and generator leakage.}
\end{solution}

\begin{solution}{iq:ml:generative-models-and-synthetic-data:5}
The generator covers only part of the distribution. Signs in returns: thin tails, too little diversity in volatility
levels, generated windows close to a few typical real ones, a two-sample classifier winning on tail features.
\iqlookfor{the definition and concrete diagnostics.}
\end{solution}

\begin{solution}{iq:ml:generative-models-and-synthetic-data:6}
Build features per window, split real windows by time blocks (with a gap longer than the window) so that no held-out
window overlaps a training one, split generated windows at random, train, report held-out accuracy with its standard error.
A random split puts near-copies on both sides and measures recognition of periods.
\iqlookfor{blocked splitting with a gap, and why overlap inflates accuracy.}
\end{solution}
